\subsubsection{Môi trường và công cụ phát triển}
\hspace*{2em} \textbf{Công nghệ:} Ứng dụng được phát triển bằng \textit{Flutter/Dart} ở phía \textit{client} và \textit{Node.js/Express} ở phía \textit{server}.  
Dữ liệu được lưu trữ trên \textit{MongoDB}, trong đó các mô hình dữ liệu (\textit{model}) bao gồm người dùng, quan hệ bạn bè, hội thoại, thành viên nhóm, tin nhắn và khóa mã hóa.  
Các thư viện mã hóa như \textit{PointyCastle}, \textit{Encrypt} và \textit{Flutter Secure Storage} được tích hợp để đảm bảo xử lý an toàn các thao tác mã hóa \textbf{RSA–AES} cũng như quản lý khóa cục bộ trên thiết bị.  
Giao tiếp thời gian thực được thực hiện thông qua \textit{WebSocket}, cho phép truyền và nhận sự kiện tức thời giữa các người dùng trong cùng một hội thoại.

\indent \textbf{Môi trường phát triển:} Ứng dụng được xây dựng và kiểm thử trên môi trường giả lập (\textit{Android Emulator}) và thiết bị thật, sử dụng \textit{Flutter CLI}, \textit{Gradle} và \textit{Android Studio} để biên dịch.  
Phía máy chủ vận hành trong môi trường \textit{Node.js} với khả năng đồng bộ dữ liệu thông qua \textit{MongoDB Change Streams}, hỗ trợ cập nhật \textit{real-time} khi có thay đổi tại tầng lưu trữ.

\indent \textbf{Cấu trúc hoạt động:} Trong giai đoạn phát triển, \textit{client} kết nối tới máy chủ thông qua địa chỉ cấu hình nội bộ.  
Khi triển khai chính thức, toàn bộ kênh truyền đều bắt buộc sử dụng \texttt{HTTPS} và \texttt{WSS}, kèm theo \textit{token} xác thực hợp lệ cho mọi giao tiếp.  
Cấu trúc này giúp hệ thống vận hành ổn định, dễ dàng mở rộng và đảm bảo an toàn dữ liệu trong suốt vòng đời hoạt động.
